% Section 13.3: Real Soft Photons; Cancellation of Divergences.
% Source: physical PDF pp. 567--571 (printed pp. 544--548).
% Starts with the section heading in the lower part of p. 567 and ends with
% Eq. (13.3.12), immediately before the Section 13.4 heading on p. 571.
% Equations: (13.3.1)--(13.3.12). Figures: none. Footnotes: two.
% Status: source-reviewed against rendered pages; no unresolved uncertainties.

\subsection{Real Soft Photons; Cancellation of Divergences}

The resolution of the infrared divergence problem encountered in the
previous section is found in the observation that it is not really possible
to measure the rate $\Gamma_{\beta\alpha}$ for a reaction $\alpha\to\beta$
involving definite numbers of photons and charged particles, because photons
of very low energy can always escape undetected. What can be measured is the
rate $\Gamma_{\beta\alpha}(E,E_T)$ for such a reaction to take place with no
unobserved photon having an energy greater than some small quantity $E$, and
with not more than some small total energy $E_T$ going into any number of
unobserved photons. (Of course, $E\leq E_T$. In an experiment without soft
photon detectors, one can rely on measurements of the energies of the `hard'
particles in $\alpha$ and $\beta$ to put a limit $E_T$ on the total energy
going into soft photons, and in this case we just set $E=E_T$.) We now turn to
a calculation of this rate.

The $S$-matrix for emitting $N$ real soft photons in a process
$\alpha\to\beta$ is obtained by contracting each of the $N$ photon
polarization indices $\mu_1,\mu_2,\ldots$ on the amplitude \eqref{eq:13.1.8}
with the appropriate coefficient function
\[
\frac{\epsilon^*_{\mu}(\mathbf q,h)}{(2\pi)^{3/2}\sqrt{2|\mathbf q|}},
\]
where $\mathbf q$ is the photon momentum, $h=\pm1$ is its helicity, and
$\epsilon^\mu$ is the corresponding photon polarization `vector'.\footnote{We
are using $\epsilon^\mu$ instead of $e^\mu$ for photon polarization vectors,
to avoid confusion with our use here of $e_n$ for electric charges.} This
gives a photon emission matrix element (the $S$-matrix element with delta
function omitted) as
\begin{align}
M^\lambda_{\beta\alpha}(\mathbf q_1,h_1,\mathbf q_2,h_2,\ldots)
={}&M^\lambda_{\beta\alpha}
\prod_{r=1}^{N}(2\pi)^{-3/2}(2|\mathbf q_r|)^{-1/2}
\sum_n\frac{\eta_n e_n[p_n\cdot\epsilon^*(\mathbf q_r,h_r)]}
{p_n\cdot q_r}.
\tag{13.3.1}\label{eq:13.3.1}
\end{align}
(The superscript $\lambda$ is to remind us that these amplitudes are to be
calculated with an infrared cutoff $\lambda$ on the momenta of virtual
photons. Eventually we shall take $\lambda\to0$. The presence of soft virtual
photons does not interfere with the result \eqref{eq:13.3.1} because of the
factorization discussed in Section 13.1.) The differential rate for emitting
$N$ soft photons into a volume $\prod_r \dd^3q_r$ of momentum space is given by
squaring this matrix element, summing over helicities, and multiplying with
$\prod_r \dd^3q_r$. We recall from Eq. \eqref{eq:8.5.7} that for $q^2=0$, the
helicity sums take the form
\begin{equation}
\sum_{h=\pm1}\epsilon^\mu(\mathbf q,h)\epsilon^{\nu*}(\mathbf q,h)
=\eta^{\mu\nu}+q^\mu c^\nu+q^\nu c^\mu,
\tag{13.3.2}\label{eq:13.3.2}
\end{equation}
where $\mathbf c\equiv-\mathbf q/2|\mathbf q|^2$ and
$c^0\equiv1/2|\mathbf q|$. The charge conservation condition
$\sum_n\eta_ne_n=0$ allows us to drop the terms in \eqref{eq:13.3.2}
involving $q_\mu$ or $q_\nu$, yielding a differential rate\footnote{The result
for $|\mathbf q|\dd\Gamma^\lambda_{\beta\alpha}(\mathbf q)/
\Gamma^\lambda_{\beta\alpha}$ in the case $N=1$ corresponds to the
distribution of energy emitted classically by a discontinuously changing
current density four-vector
$J^\mu(x)=\sum_n\eta_ne_n\delta^3(\mathbf x-\mathbf v_nt)p_n^\mu/E_n$,
with the sum here running only over particles in the initial state for $t<0$
and in the final state for $t>0$.}
\begin{align}
\dd\Gamma^\lambda_{\beta\alpha}(\mathbf q_1,\ldots,\mathbf q_N)
={}&\Gamma^\lambda_{\beta\alpha}
\prod_{r=1}^{N}\frac{\dd^3q_r}{(2\pi)^3(2|\mathbf q_r|)}
\sum_{nm}\frac{\eta_n\eta_m e_ne_m(p_n\cdot p_m)}
{(p_n\cdot q_r)(p_m\cdot q_r)}.
\tag{13.3.3}\label{eq:13.3.3}
\end{align}

To calculate the differential rate for the emission of $N$ soft photons with
definite energies $\omega_r=|\mathbf q_r|$ we must integrate
\eqref{eq:13.3.3} over the directions of the photon momenta $\mathbf q_r$.
These integrals are the same as those we encountered in the integrals
\eqref{eq:13.2.8},
\begin{equation}
-\pi(p_n\cdot p_m)\int
\frac{\dd^2\hat{\mathbf q}}
{(E_n-\hat{\mathbf q}\cdot\mathbf p_n)
 (E_m-\hat{\mathbf q}\cdot\mathbf p_m)}
=\frac{2\pi^2}{\beta_{nm}}
\ln\left(\frac{1+\beta_{nm}}{1-\beta_{nm}}\right).
\tag{13.3.4}\label{eq:13.3.4}
\end{equation}
Integrating \eqref{eq:13.3.3} over photon directions thus gives the
differential rate for photons of energy $\omega_1,\ldots,\omega_N$:
\begin{equation}
\dd\Gamma^\lambda_{\beta\alpha}(\omega_1,\ldots,\omega_N)
=\Gamma^\lambda_{\beta\alpha}A(\alpha\to\beta)^N
\frac{\dd\omega_1}{\omega_1}\cdots\frac{\dd\omega_N}{\omega_N},
\tag{13.3.5}\label{eq:13.3.5}
\end{equation}
where $A(\alpha\to\beta)$ is the same constant encountered in the previous
section:
\[
A(\alpha\to\beta)=-\frac{1}{8\pi^2}\sum_{nm}
\frac{e_ne_m\eta_n\eta_m}{\beta_{nm}}
\ln\left(\frac{1+\beta_{nm}}{1-\beta_{nm}}\right).
\]

We see from \eqref{eq:13.3.5} that an unrestricted integral over the energies
of the emitted photons would introduce another infrared divergence. However,
unitarity demands that if we use an infrared cutoff for the momenta of the
virtual photons (as implied by the superscript $\lambda$) then we must use the
same infrared cutoff for the real photons. To calculate the rate
$\Gamma^\lambda_{\beta\alpha}(E,E_T)$ for the reaction $\alpha\to\beta$ with
not more than energy $E$ going into any one unobserved photon and not more than
energy $E_T$ going into any number of unobserved photons (with $E$ and $E_T$
chosen small enough to justify the approximations used in deriving
\eqref{eq:13.3.1}), we must integrate \eqref{eq:13.3.5} over all photon
energies, subject to the limits $E\geq\omega_r\geq\lambda$ and
$\sum_r\omega_r\leq E_T$, then divide by $N!$ because this integral includes
configurations that differ only by permutations of the $N$ soft photons, and
finally sum over $N$. This gives
\begin{equation}
\Gamma^\lambda_{\beta\alpha}(E,E_T)
=\Gamma^\lambda_{\beta\alpha}\sum_{N=0}^{\infty}
\frac{A(\alpha\to\beta)^N}{N!}
\int_{\substack{E\geq\omega_r\geq\lambda\\
                 \sum_r\omega_r\leq E_T}}
\prod_{r=1}^{N}\frac{\dd\omega_r}{\omega_r}.
\tag{13.3.6}\label{eq:13.3.6}
\end{equation}
This integral would factor into the product of $N$ integrals over the
individual $\omega_r$ were it not for the restriction
$\sum_r\omega_r\leq E_T$. This restriction may be implemented by including as
a factor in the integrand a step function
\begin{equation}
\theta\left(E_T-\sum_r\omega_r\right)
=\frac{1}{\pi}\int_{-\infty}^{\infty}\dd u\,
\frac{\sin E_Tu}{u}\exp\left(iu\sum_r\omega_r\right).
\tag{13.3.7}\label{eq:13.3.7}
\end{equation}
Eq. \eqref{eq:13.3.6} then becomes
\begin{equation}
\Gamma^\lambda_{\beta\alpha}(E,E_T)
=\frac{1}{\pi}\int_{-\infty}^{\infty}\dd u\,
\frac{\sin E_Tu}{u}
\exp\left(A(\alpha\to\beta)\int_\lambda^E
\frac{\dd\omega}{\omega}e^{\ii\omega u}\right)
\Gamma^\lambda_{\beta\alpha}.
\tag{13.3.8}\label{eq:13.3.8}
\end{equation}
The integral in the exponent can be done in the limit $\lambda\ll E_T$ by
writing it as the sum of the integral of $(e^{\ii\omega u}-1)/\omega$, in which
we can set $\lambda=0$, and the integral of $1/\omega$, which is trivial.
Rescaling the $u$ and $\omega$ variables, this gives for $\lambda\ll E$:
\begin{equation}
\Gamma^\lambda_{\beta\alpha}(E,E_T)\to
\mathcal F\bigl(E/E_T;A(\alpha\to\beta)\bigr)
\left(\frac{E}{\lambda}\right)^{A(\alpha\to\beta)}
\Gamma^\lambda_{\beta\alpha},
\tag{13.3.9}\label{eq:13.3.9}
\end{equation}
where
\begin{align}
\mathcal F(x;A)
\equiv{}&\frac{1}{\pi}\int_{-\infty}^{\infty}\dd u\,
\frac{\sin u}{u}\exp\left(A\int_0^x\frac{\dd\omega}{\omega}
(e^{\ii\omega u}-1)\right)\notag\\
={}&1-\frac{A^2\theta(x-\tfrac12)}{2}
\int_{1-x}^{x}\frac{\dd\omega}{\omega}
\ln\left(\frac{x}{1-\omega}\right)+\cdots.
\tag{13.3.10}\label{eq:13.3.10}
\end{align}
For $E$ and $E_T$ of the same order and $A\ll1$ the factor
$\mathcal F(E/E_T;A)$ in \eqref{eq:13.3.9} is close to unity; for instance,
\[
\mathcal F(1;A)\simeq1-\frac{1}{12}\pi^2A^2+\cdots.
\]

Because $A(\alpha\to\beta)>0$, the factor
$(E/\lambda)^{A(\alpha\to\beta)}$ in \eqref{eq:13.3.9} becomes infinite in
the limit $\lambda\to0$. However, Eq. \eqref{eq:13.2.10} shows that the rate
$\Gamma^\lambda_{\beta\alpha}$ vanishes in this limit:
\[
\Gamma^\lambda_{\beta\alpha}
=\left(\frac{\lambda}{\Lambda}\right)^{A(\alpha\to\beta)}
\Gamma^\Lambda_{\beta\alpha}.
\]
Using this in \eqref{eq:13.3.9} shows that the infrared cutoff $\lambda$ drops
out in the limit $\lambda\ll E$:
\begin{equation}
\Gamma^\lambda_{\beta\alpha}(E,E_T)\to
\mathcal F\bigl(E/E_T;A(\alpha\to\beta)\bigr)
\left(\frac{E}{\Lambda}\right)^{A(\alpha\to\beta)}
\Gamma^\Lambda_{\beta\alpha}.
\tag{13.3.11}\label{eq:13.3.11}
\end{equation}

We remind the reader that the energy $\Lambda$ is just a convenient dividing
point between `soft' photons which are taken into account explicitly in
\eqref{eq:13.3.11} and `hard' photons whose effects are buried in
$\Gamma^\Lambda_{\beta\alpha}$. The right-hand side of \eqref{eq:13.3.11} is
independent of $\Lambda$ because
$\Gamma^\Lambda_{\beta\alpha}\propto\Lambda^A$. However, in theories with a
small coupling constant like quantum electrodynamics it is frequently a good
strategy to take $\Lambda$ to be sufficiently small compared with the typical
energies $W$ involved in the collision so that the approximations made here
apply for photon energies less than $\Lambda$, but large enough so that
$A(\alpha\to\beta)\ln(W/\Lambda)\ll1$. Then it may be a good approximation to
calculate $\Gamma^\Lambda_{\beta\alpha}$ in lowest-order perturbation theory,
with the dominant radiative corrections for $E\ll\Lambda$ given by the factor
$(E/\Lambda)^A$ in \eqref{eq:13.3.11}.

\begin{center}
$*\ *\ *$
\end{center}

The same cancellation of infrared divergences occurs for soft
gravitons~\hyperref[ref:13.3]{[3]}. The rate for any process
$\alpha\to\beta$, with not more than an energy $E$ going into soft gravitons,
turns out to be proportional to $E^B$, where
\begin{equation}
B=\frac{G}{2\pi}\sum_{nm}\eta_n\eta_m m_nm_m
\frac{1+\beta_{nm}^2}{\beta_{nm}\sqrt{1-\beta_{nm}^2}}
\ln\left(\frac{1+\beta_{nm}}{1-\beta_{nm}}\right).
\tag{13.3.12}\label{eq:13.3.12}
\end{equation}
